% TOP_PROBLEM: {"id": 3208, "title": "Mapping planar graphs to odd cycles", "category_id": 3, "category": {"id": 3, "name": "graph_theory", "display_name": "Graph Theory", "description": "Problems involving graphs, networks, and their properties.", "slug": "graph-theory", "order_index": 3, "created_at": "2026-07-31T15:26:25.671Z"}, "status": "open", "problem_number": "OPG-412", "set_id": 12, "statusQualification": "Ordinary girth at least 4k is preserved from the source; no unsupported equivalence with an odd-girth-only hypothesis is asserted."}
% FIRST_POSTED: 2026-10-01
% ATTEMPT_STATUS: unresolved
% MODEL: GPT 6 Astra Ultra
% UnsolvedMath ID 3208; source catalogue code OPG-412
% !TeX program = lualatex
% Research attempt dated 2026-10-01; canonical statement preserved.
\documentclass[11pt,a4paper]{article}
\usepackage[margin=25mm]{geometry}
\usepackage{fontspec}
\setmainfont{lmroman10-regular.otf}[BoldFont=lmroman10-bold.otf,
 ItalicFont=lmroman10-italic.otf,BoldItalicFont=lmroman10-bolditalic.otf]
\usepackage{amsmath,amssymb,mathtools}
\usepackage{unicode-math}
\setmathfont{latinmodern-math.otf}
\AtBeginDocument{\let\setminus\smallsetminus}
\usepackage{xurl,hyperref}
\hypersetup{colorlinks=true,urlcolor=blue}
\setcounter{secnumdepth}{0}
\setlength{\parindent}{0pt}
\setlength{\parskip}{5pt}
\setlength{\emergencystretch}{3em}
\begin{document}
\section*{Mapping planar graphs to odd cycles}\label{3208}
{\small UnsolvedMath ID 3208 · OPG-412 · Graph Theory · OpenGarden}
\subsection{Definitions and mathematical statement}
For an integer $k\geq1$, let $C_{2k+1}$
be the odd cycle on residues
$\mathbb Z/(2k+1)\mathbb Z$, with adjacent
residues differing by $1$ or $-1$.
A homomorphism from a simple graph $G$
to this cycle is a function
$f:V(G)\to\mathbb Z/(2k+1)\mathbb Z$
such that
\[
 f(v)-f(u)\equiv\pm1\pmod{2k+1}
 \qquad(uv\in E(G)).
\]
There is no injectivity or surjectivity
requirement on $f$.

A finite graph is planar if it has a
drawing in the plane without edge
crossings. Its girth is the length of
its shortest simple cycle, with girth
$+\infty$ for a forest. The source
conjecture is the assertion
\[
 \forall k\geq1\quad\forall G\text{ finite simple planar}:
 \operatorname{girth}(G)\geq4k
 \Longrightarrow G\to C_{2k+1}.
\]
Thus the cycle target becomes longer
as the required girth increases.

At $k=1$ the target is a triangle and
the desired map is exactly a proper
three-colouring. For larger $k$, merely
using $2k+1$ distinct labels is not
enough: the allowed label pairs are
only edges of the target cycle.
The hypothesis here concerns ordinary
girth, excluding short even cycles
as well as odd ones. This preserves
the supplied catalogue formulation;
an odd-girth-only hypothesis is not
silently substituted. The statement
includes disconnected graphs and
forests under the conventions above.
\subsection{Short English statement}
Must every planar graph with no cycle shorter than 4k map homomorphically to the odd cycle of length 2k+1?
\subsection{Research attempt}
\paragraph{Scope and method.}
This research attempt by GPT-6 Astra Ultra is dated 1 October 2026.
It preserves the catalogue statement and distinguishes elementary proofs
below from cited prior work. No novelty claim or formal proof-assistant
verification is made. The final paragraph states the precise remaining scope.
\paragraph{Approach and source qualification.}
The source [S2] attributes the conjecture to Jaeger and describes the
subdivided-wheel obstruction below. We separate that sharpness example
from an elementary positive extension method. Throughout, girth means
the shortest cycle of either parity; replacing it by odd girth changes
the problem. Put $q=2k+1$ and identify the target vertices with
$\mathbb Z/q\mathbb Z$.

\paragraph{An exact path-extension lemma.}
If a path has length $L$, assigning its edge steps to $\{1,-1\}$ gives
endpoint differences $L-2j$ modulo $q$, where $0\le j\le L$.
When $L\ge2k=q-1$, the first $q$ of these values represent all residues:
if $L-2i\equiv L-2j\pmod q$, oddness of $q$ implies $i\equiv j\pmod q$.
Therefore every assignment of endpoint colours extends over a path of
length at least $2k$. This is a direct existence proof: choose the required
number of negative steps and place them in any order. No injectivity is
required of a homomorphism.

At length $L=2k-1=q-2$, exactly $q-1$ residues occur and the missing
residue is zero. Indeed a zero would be an odd multiple of $q$ with
absolute value at most $q-2$, which is impossible. Thus identifying the
endpoint colours across such a path can prevent extension. The one-edge
difference in this threshold matters in any reduction that suppresses
degree-two vertices.

\paragraph{The lower-bound obstruction.}
Take a rim $C_{4k-1}$ and a hub, replacing each spoke by an internally
disjoint path of length $2k-1$. This is planar. A cycle avoiding the hub
is the rim; a cycle using the hub contains two spokes and at least one
rim edge. Its girth is consequently $4k-1$. In a hypothetical target
mapping with hub colour $a$, no rim vertex could have colour $a$, by the
preceding missing-residue calculation. The rim would then map to
$C_{2k+1}-a$, a path and hence a bipartite graph. This cannot receive a
homomorphism from an odd cycle. Thus girth $4k-1$ is insufficient. For
$k=1$ the construction is simply $K_4$, which does not map to $C_3$.
This argument verifies sharpness below the proposed threshold; it is
not a counterexample at the requested girth $4k$.

\paragraph{A positive theorem for cactus graphs.}
An even cycle maps onto any chosen edge of the target by alternating its
endpoints. An odd cycle of length $m\ge q$ maps by traversing the target
once and inserting $(m-q)/2$ backtracking pairs, giving a closed walk of
length $m$. Rotation permits any desired image of one fixed vertex.
In a cactus, every block is an edge or a cycle. Root the block-cutvertex
tree and apply these constructions successively, agreeing with the
already assigned image of the shared cutvertex. This proves that every
cactus all of whose odd cycles have length at least $2k+1$ maps to
$C_{2k+1}$. In particular the catalogue's conclusion holds for every
cactus of girth at least $4k$, without needing a planar embedding.

\paragraph{Attempted induction and its failure point.}
In a vertex-minimal counterexample one can remove a leaf and later extend
its image. One can also remove the internal vertices of a sufficiently
long degree-two path and extend any mapping of the remaining graph by
the path lemma. But a general planar graph of the stated girth need not
offer this precise removable path at each induction step. Suppressing a
shorter thread can change girth and can force endpoint colours that fail
the missing-residue test. The cactus argument avoids this issue because
blocks meet in only one vertex. General planar cycles can overlap along
paths or interact through several prescribed endpoints.

\paragraph{Verification and final status.}
The root batch script computes all walk residues for $1\le k\le8$ and
checks both path thresholds exactly. The symbolic arguments above cover
all $k$, including the triangle target. \textbf{UNRESOLVED.} We verify the
sharpness obstruction and a positive cactus theorem. A reduction preserving
the $4k$ girth hypothesis and all necessary endpoint constraints is still
required for arbitrary planar graphs.

\subsection{Sources}
\begin{itemize}
\item[S1] ULAM.AI and the UnsolvedMath Contributors, supplied problems.json, record 3208 (OPG-412), OpenGarden collection. Catalogue identity and source transcription; CC BY 4.0.
\url{https://huggingface.co/datasets/ulamai/UnsolvedMath}.
\item[S2] Open Problem Garden, Mapping planar graphs to odd cycles. Original problem and discussion; the mathematical formulation below is expanded from the supplied source record.
\url{http://www.openproblemgarden.org/op/mapping_planar_graphs_to_odd_cycles}.

\end{itemize}
\end{document}
